\documentclass{article}
\usepackage[UTF8]{ctex}
\usepackage{amsmath,amssymb,amsfonts}
\usepackage{enumitem}
\usepackage{geometry}
\geometry{a4paper, margin=1in}
% === 错题本排版优化（不参与题目标识） ===
\setlength{\parindent}{0pt}
\setlength{\parskip}{0.35em plus 0.10em minus 0.08em}
\setlength{\emergencystretch}{3em}
\setlist{itemsep=0.25em, topsep=0.40em, parsep=0pt}
\allowdisplaybreaks[2]
% === 错题本排版优化结束 ===
\begin{document}
\section*{北京交通大学2025年高等代数真题}
\section*{5．北京交通大学2025年高等代数试卷（回忆版）}
一．填空题。\\

\subsection{题目1}
1．在有理数域上分解多项式 \(f(x)=x^{4}-3 x^{3}+6 x-4\) \(\_\_\_\_\)．

\subsection{题目2}
2．设 \(B=\left(\begin{array}{lll}3 & 0 & 1 \\ 1 & 1 & 0 \\ 1 & 0 & 1\end{array}\right), A=\left(\alpha_{1}, \alpha_{2}, \alpha_{3}\right)\) 是 3 阶方阵且 \(A^{*}\) 是非零矩阵，且 \(\alpha_{1}+\alpha_{2}-2 \alpha_{3}=0\) ，求

\subsection{题目3}
3．设 \(A=\left(\begin{array}{ccccc}0 & a_{1} & 0 & \cdots & 0 \\ 0 & 0 & a_{2} & \cdots & 0 \\ \vdots & \vdots & \vdots & & \vdots \\ 0 & 0 & 0 & \cdots & a_{n-1} \\ a_{n} & 0 & 0 & \cdots & 0\end{array}\right)\) ，求 \(A^{-1}=\) \(\_\_\_\_\)，其中 \(a_{i} \neq 0(i=1,2, \cdots, n)\) ．\\
\(\_\_\_\_\)

\subsection{题目4}
4．线性方程组

\[
\left\{\begin{array}{l}
x_{1}+x_{2}+x_{3}+x_{4}=1 \\
4 x_{1}+3 x_{2}+5 x_{3}-x_{4}=-1 \\
a x_{1}+x_{2}+3 x_{3}+b x_{4}=1
\end{array}\right.
\]

有 3 个线性无关的解，则 \(a, b\) 的值为 \(\_\_\_\_\)．

\subsection{题目5}
5．已知行列式 \(|A|\) 的每行元素之和为 \(a\) ，且 \(|A|=b\) ，则 \(A_{1 n}+A_{2 n}+\cdots+A_{n n}=\) \(\_\_\_\_\)．

\subsection{题目6}
6．\(X^{\mathrm{T}} A X\) 的符号差为 1 ，且 \(A^{2}-2 A=3 E\) ，则 \(|A|=\) \(\_\_\_\_\)．

\subsection{题目7}
7．若线性变换 \(\sigma\) 有 \(n\) 个不同的特征值，则 \(\sigma\) 有 \(\_\_\_\_\)个不变子空间。

\subsection{题目8}
8．\(A=\left(\begin{array}{ccc}3 & -2 & 0 \\ -2 & 2 & -2 \\ 0 & -2 & 1\end{array}\right), A+t E\) 正定，\(t\) 满足 \(\_\_\_\_\)．

\subsection{题目9}
9．\(A=\left(\begin{array}{ccc}3 & 0 & 8 \\ 3 & -1 & 6 \\ -2 & 0 & -5\end{array}\right)\) 的最小多项式为 \(\_\_\_\_\)，Jordan 标准型为 \(\_\_\_\_\)．

\subsection{题目10}
10．设 \((\alpha, \beta)=\alpha^{\mathrm{T}} A \beta\) ，其中 \(A=\left(\begin{array}{ll}5 & 2 \\ 2 & 1\end{array}\right)\) ，则 \(\binom{1}{0},\binom{1}{1}\) 的度量阵为 \(\_\_\_\_\)．

二．计算行列式

\[
D=\left|\begin{array}{ccccc}
x & a_{2} & a_{3} & \cdots & a_{n} \\
a_{1} & x & a_{3} & \cdots & a_{n} \\
a_{1} & a_{2} & x & \cdots & a_{n} \\
\vdots & \vdots & \vdots & & \vdots \\
a_{1} & a_{2} & a_{3} & \cdots & x
\end{array}\right|, \text { 其中 } x \neq a_{i}(i=1,2, \cdots, n) \text {. }
\]

三．已知非齐次线性方程组

\[
\left\{\begin{array}{l}
x_{1}+x_{2}-2 x_{3}+3 x_{4}=0 \\
2 x_{1}+x_{2}-6 x_{3}+4 x_{4}=-1 \\
3 x_{1}+2 x_{2}+p x_{3}+7 x_{4}=-1 \\
x_{1}-x_{2}-6 x_{3}-x_{4}=t
\end{array}\right.
\]

有无穷多解，求 \(p, t\) 的值与方程组的通解．

四．设 \(A=\left(\begin{array}{ccc}0 & a & a^{2} \\ \frac{1}{a} & 0 & a \\ \frac{1}{a^{2}} & \frac{1}{a} & 0\end{array}\right)\) ，求三阶可逆矩阵 \(P\) ，使得 \(P^{-1} A P=\Lambda\) ，且求 \(A^{n}\) ．

五．已知

\[
\alpha_{1}=(1,1, a)^{\mathrm{T}}, \alpha_{2}=(1, a, 1)^{\mathrm{T}}, \alpha_{3}=(a, 1,1)^{\mathrm{T}}, \beta=(1,1,-2)^{\mathrm{T}} .
\]

\(\beta\) 可由 \(\alpha_{1}, \alpha_{2}, \alpha_{3}\) 线性表出且不唯一，求 \(a\) 的值。记 \(A=\left(\alpha_{1}, \alpha_{2}, \alpha_{3}\right)\) ，求正交变换 \(Q\) ，通过 \(X=Q Y\)，使

\[
f\left(x_{1}, x_{2}, x_{3}\right)=X^{\mathrm{T}}\left(E_{3}+A\right)^{-1} X
\]

为标准型。

六。设 \(f(x)=\left(x-k_{1}\right)\left(x-k_{2}\right) \cdots\left(x-k_{n}\right)+1\) ，其中 \(k_{1}, k_{2}, \cdots, k_{n}(n>2)\) 是互异的整数。证明：\(f(x)\)在有理数域上可约的充要条件是 \(f(x)\) 是某个整系数多项式的完全平方．

\begin{itemize}
  \item[七．] \(V_{1}=L\left(\alpha_{1}, \alpha_{2}, \alpha_{3}\right), V_{2}=L\left(\beta_{1}, \beta_{2}\right)\) ．（具体数据忘了）
  \item[（1）] 求 \(V_{1}+V_{2}\) 的维数和一组基。
  \item[（2）] 求 \(V_{1} \cap V_{2}\) 的维数和一组基。
  \item[（3）] 证明：．．．
\end{itemize}

八．设 \(\sigma, \tau\) 为 \(n\) 维欧氏线性空间 \(V\) 上的线性变换，且满足 \((\sigma(\alpha), \beta)=(\alpha, \tau(\beta))\) 对任意的 \(\alpha, \beta \in V\) ．证明：

\begin{itemize}
  \item[（1）] \(\operatorname{Ker} \sigma=(\operatorname{Im} \tau)^{\perp}\) ．
  \item[（2）] \(V=\operatorname{Ker} \sigma \oplus \operatorname{Im} \tau\) ．
\end{itemize}

九．已知矩阵 \(B\) 为正定矩阵，\(A-B\) 为半正定矩阵，证明：

\begin{itemize}
  \item[（1）] \(|A-\lambda B|=0\) 的根均大于等于1 。
  \item[（2）] \(|A| \geq|B|\) ．
\end{itemize}

十．也有同学回顾共十个大题，缺少的题目未知…
\end{document}
